\documentclass[12pt,a4paper]{article}
\usepackage[margin=0.6in]{geometry} % പേജ് മാർജിൻ ക്രമീകരിക്കുന്നു
\usepackage{fontspec}
\usepackage{array}

% മലയാളം ഫോണ്ട് ക്രമീകരണം (സിസ്റ്റത്തിലുള്ള Rachana അല്ലെങ്കിൽ Nirmala UI ഉപയോഗിക്കുക)
\setmainfont[Script=Malayalam]{Rachana} 

\pagestyle{empty} % പേജ് നമ്പർ ഒഴിവാക്കാൻ

\sloppy 
\emergencystretch=3em 

\begin{document}
	
	% മുകൾഭാഗം: വിലാസവും സ്റ്റാമ്പ് പതിക്കാനുള്ള ബോക്സും സൈഡ്-ബൈ-സൈഡ് ആയി ക്രമീകരിക്കുന്നു
	\noindent
	\begin{minipage}[t]{0.58\textwidth}
		\textbf{അയക്കുന്നയാൾ,} \\[0.1cm]
		പേര്: \dotfill\null \\[0.15cm]
		വീട്ടുപേര് / വീട്ടുവിലാസം: \dotfill\null \\[0.15cm]
		\hspace*{3.4cm} \dotfill\null \\[0.15cm]
		സ്ഥലം: \dotfill\null \\[0.15cm]
		പോസ്റ്റ് ഓഫീസ് / ജില്ല: \dotfill\null \\[0.15cm]
		മൊബൈൽ നമ്പർ: \dotfill\null
	\end{minipage}
	\hfill
	\begin{minipage}[t]{0.38\textwidth}
		\centering
		\setlength{\fboxrule}{0.8pt}
		\framebox[\textwidth]{\parbox[c][3.2cm][c]{\textwidth}{\centering
				\small
				\textbf{Rs.10/- കോടതി} \\
				\textbf{ഫീ സ്റ്റാമ്പ്} \\
				\textbf{ഇവിടെ പതിക്കുക} \\
				\vspace{0.2cm}
				(RTI അപേക്ഷ ഫീസ്)
		}}
	\end{minipage}
	
	\vspace{0.4cm}
	
	\noindent
	\textbf{സ്വീകർത്താവ്,} \\
	പബ്ലിക് ഇൻഫർമേഷൻ ഓഫീസർ, \\
	\makebox[6cm]{\dotfill} ഗ്രാമപഞ്ചായത്ത് ഓഫീസ്, \\
	\makebox[6cm]{\dotfill} ജില്ല.
	
	\vspace{0.3cm}
	
	\noindent
	\textbf{വിഷയം:} വിവരാവകാശ നിയമം – 2005 പ്രകാരം ഭിന്നശേഷി കുട്ടികളുടെ സ്കോളർഷിപ്പ് സംബന്ധിച്ച വിവരങ്ങൾ ലഭ്യമാക്കുന്നത്.
	
	\vspace{0.2cm}
	
	\noindent
	\textbf{സർ,} \\
	ഞാൻ \makebox[3.0cm]{\dotfill} വാർഡിൽ താമസിക്കുന്ന \makebox[6cm]{\dotfill} എന്ന ഭിന്നശേഷിക്കാരനായ/യായ കുട്ടിയുടെ രക്ഷിതാവാണ് (അച്ഛൻ / അമ്മ / മറ്റ്). എന്റെ മകന്/മകൾക്ക് വർഷംതോറും ലഭിക്കാറുള്ള ഡിസിബിലിറ്റി സ്കോളർഷിപ്പിനായി \makebox[2.5cm]{\dotfill} സാമ്പത്തിക വർഷത്തിൽ \makebox[3.5cm]{\dotfill} വഴി (അംഗനവാടി / പഞ്ചായത്ത് / മറ്റ്) അപേക്ഷ നൽകിയിരുന്നു. എന്നാൽ അനുവദിക്കേണ്ട തുകയിൽ നിന്നും Rs. \makebox[2.5cm]{\dotfill}/- രൂപ മാത്രമാണ് ലഭിച്ചത്. ബാക്കി തുകയായ Rs. \makebox[2.5cm]{\dotfill}/- രൂപ ഇതുവരെയും ലഭിച്ചിട്ടില്ല. ആയതിനാൽ താഴെ പറയുന്ന ചോദ്യങ്ങൾക്കുള്ള മറുപടി വിവരാവകാശ നിയമം – 2005 പ്രകാരം ലഭ്യമാക്കാൻ താത്പര്യപ്പെടുന്നു.
	
	\vspace{0.2cm}
	
	\noindent
	\textbf{ചോദ്യങ്ങൾ:}
	\begin{enumerate}
		\setlength{\itemsep}{1pt}
		\setlength{\parskip}{0pt}
		\setlength{\parsep}{0pt}
		\setlength{\topsep}{0pt}
		\setlength{\partopsep}{0pt}
		\item \makebox[6cm]{\dotfill} ഗ്രാമപഞ്ചായത്തിൽ ആകെ എത്ര ഭിന്നശേഷി കുട്ടികൾക്കാണ് സ്കോളർഷിപ്പ് അനുവദിക്കുന്നത്?
		\item ഇതിൽ എത്ര കുട്ടികൾക്ക് പൂർണ്ണമായ സ്കോളർഷിപ്പ് തുക അനുവദിച്ചിട്ടുണ്ട്?
		\item എത്ര കുട്ടികൾക്കാണ് ഇനി സ്കോളർഷിപ്പ് തുക ലഭിക്കാൻ ബാക്കിയുള്ളത്?
		\item സ്കോളർഷിപ്പ് തുക ലഭിച്ചവർക്ക് എത്ര രൂപ വീതമാണ് നൽകിയിട്ടുള്ളത്?
		\item കഴിഞ്ഞ സാമ്പത്തിക വർഷം (\makebox[2cm]{\dotfill}) ഇതിനായി ബജറ്റിൽ വകയിരുത്തിയ തുക എത്ര?
		\item ബജറ്റിൽ ഉൾപ്പെടുത്തിയ തുകയിൽ നിന്നും കഴിഞ്ഞ സാമ്പത്തിക വർഷത്തിൽ ഇതിനായി എത്ര രൂപ ചെലവഴിച്ചു?
		\item ഭിന്നശേഷി കുട്ടികൾക്ക് പൂർണ്ണ തുക നൽകണമെന്ന സർക്കാർ ഉത്തരവുണ്ടായിട്ടും എന്തുകൊണ്ട് പൂർണ്ണ തുക അനുവദിച്ചില്ല?
		\item കഴിഞ്ഞ സാമ്പത്തിക വർഷത്തെ ഭിന്നശേഷി കുട്ടികൾക്ക് നൽകാനുള്ള ബാക്കി സ്കോളർഷിപ്പ് തുക എപ്പോൾ ലഭ്യമാക്കും?
		\item ബ്ലോക്ക് പഞ്ചായത്തും ജില്ലാ പഞ്ചായത്തും അവരുടെ വിഹിതം കൃത്യമായി നൽകിയിട്ടുണ്ടോ? എങ്കിൽ:
		\begin{enumerate}
			\setlength{\itemsep}{0pt}
			\setlength{\parskip}{0pt}
			\setlength{\parsep}{0pt}
			\setlength{\topsep}{0pt}
			\item[a)] ഓരോരുത്തരും എത്ര രൂപ വീതം നൽകണമായിരുന്നു?
			\item[b)] യഥാർത്ഥത്തിൽ എത്ര രൂപയാണ് നൽകിയിട്ടുള്ളത്?
		\end{enumerate}
	\end{enumerate}
	
	\vspace{0.6cm}
	
	\noindent
	\begin{minipage}[t]{0.45\textwidth}
		സ്ഥലം: \dotfill\null \\[0.15cm]
		തീയതി: \dotfill\null
	\end{minipage}
	\hfill
	\begin{minipage}[t]{0.45\textwidth}
		\raggedleft
		\textbf{വിശ്വസ്തതയോടെ,} \\[0.5cm]
		ഒപ്പ്: \dotfill\null \\[0.15cm]
		പേര്: \dotfill\null
	\end{minipage}
	
\end{document}